% part: intuitionistic-logic
% chapter: introduction
% section: constructive-reasoning

\documentclass[../../../include/open-logic-section]{subfiles}

\begin{document}

\olfileid{int}{int}{cr}

\olsection{في الاستدلال البنائي}

ليس المنطق الحدسي توسعة للمنطق الكلاسيكي، كتوسيعه بالمؤثرات الجهوية أو
بكمّيات الرتبة الثانية، بل هو تقييد له؛ ومن هذه الجهة يسمى «غير كلاسيكي».
وذلك أن الاستدلال الكلاسيكي \emph{غير بنائي} من وجوه، وأما المنطق الحدسي
فيراد به ضبط استدلال أشد بناءً، تقوم عليه طائفة من الرياضيات البنائية.
وتكشف الأمثلة الآتية عن موجب هذا التقييد.

لندّع أن رجلًا زعم أنه عيّن عددًا طبيعيًا~${\OLId{latin-ordinary}{n}}$ بهذه الصفة: إن كان ${\OLId{latin-ordinary}{n}}$
زوجيًا فحدسية ريمان صحيحة، وإن كان ${\OLId{latin-ordinary}{n}}$ فرديًا فهي خاطئة. وهذا في ظاهره
خبر عظيم!{} فصدق حدسية ريمان أو كذبها من كبار المسائل المفتوحة، وقد صار
الأمر --- فيما يزعم --- حسابًا يسيرًا: أننظر أزوجي ذلك العدد أم فردي.

فما ذلك العدد~${\OLId{latin-ordinary}{n}}$؟ يقول: ${\OLId{latin-ordinary}{n}}$ هو العدد الطبيعي الذي يساوي ${\OLMathNumeral{2}}$ إن صحت
حدسية ريمان، ويساوي ${\OLMathNumeral{3}}$ إن لم تصح.

فلك أن تغضب وتطلب رد مالك. نعم، يعيّن هذا الوصف عند الكلاسيكي قيمة واحدة
لـ${\OLId{latin-ordinary}{n}}$؛ ولكن المطلوب قيمة لـ${\OLId{latin-ordinary}{n}}$ معطاة \emph{صراحة}، لا وصفًا يتوقف على
حل المسألة نفسها.

ومثال آخر أقرب إلى صناعة الرياضيات: نعلم أن قوة عدد غير نسبي قد تكون
نسبية إذا كان الأس نسبيًا؛ فمثلًا~$\sqrt{{\OLMathNumeral{2}}}^{\OLMathNumeral{2}} = {\OLMathNumeral{2}}$. فهل يقع مثل ذلك إذا
كان الأس نفسه \emph{غير نسبي}؟ تجيب المبرهنة بالإيجاب:

\begin{thm}
يوجد عددان غير نسبيين ${\OLId{latin-ordinary}{a}}$ و${\OLId{latin-ordinary}{b}}$ بحيث يكون ${\OLId{latin-ordinary}{a}}^{\OLId{latin-ordinary}{b}}$ نسبيًا.
\end{thm}

\begin{proof}
انظر في $\sqrt{{\OLMathNumeral{2}}}^{\sqrt{{\OLMathNumeral{2}}}}$. إذا كان هذا العدد نسبيًا فقد انتهينا:
نأخذ ${\OLId{latin-ordinary}{a}} = {\OLId{latin-ordinary}{b}} = \sqrt{{\OLMathNumeral{2}}}$. وإلا فهو غير نسبي. وعندئذ
\[
(\sqrt{{\OLMathNumeral{2}}}^{\sqrt{{\OLMathNumeral{2}}}})^{\sqrt{{\OLMathNumeral{2}}}} = \sqrt{{\OLMathNumeral{2}}}^{\sqrt{{\OLMathNumeral{2}}} \cdot
  \sqrt{{\OLMathNumeral{2}}}} = \sqrt{{\OLMathNumeral{2}}}^{\OLMathNumeral{2}} = {\OLMathNumeral{2}},
\]
وهو عدد نسبي. ولذلك نأخذ في هذه الحالة
${\OLId{latin-ordinary}{a}}=\sqrt{{\OLMathNumeral{2}}}^{\sqrt{{\OLMathNumeral{2}}}}$ و${\OLId{latin-ordinary}{b}}=\sqrt {\OLMathNumeral{2}}$.
\end{proof}

وهذا برهان صحيح عند أكثر الرياضيين؛ غير أنه لا يعيّن الزوج الذي أثبت
وجوده، وفي ذلك ما لا يرضي غرض البناء. ويمكن إثبات النتيجة نفسها مع إعطاء
الزوج ${\OLId{latin-ordinary}{a}}$، ${\OLId{latin-ordinary}{b}}$ \emph{في} البرهان: خذ
${\OLId{latin-ordinary}{a}} = \sqrt{{\OLMathNumeral{3}}}$ و${\OLId{latin-ordinary}{b}} = \log_{\OLMathNumeral{3}} {\OLMathNumeral{4}}$. فالعدد~$\sqrt{\OLMathNumeral{3}}$ غير نسبي، وكذلك
$\log_{\OLMathNumeral{3}} {\OLMathNumeral{4}}$ غير نسبي؛ إذ لو كان $\log_{\OLMathNumeral{3}} {\OLMathNumeral{4}}={\OLId{latin-ordinary}{m}}/{\OLId{latin-ordinary}{n}}$ لعددين طبيعيين موجبين
${\OLId{latin-ordinary}{m}},{\OLId{latin-ordinary}{n}}$، لكان ${\OLMathNumeral{3}}^{\OLId{latin-ordinary}{m}}={\OLMathNumeral{4}}^{\OLId{latin-ordinary}{n}}={\OLMathNumeral{2}}^{{\OLMathNumeral{2}}{\OLId{latin-ordinary}{n}}}$، وهذا يناقض وحدانية التحليل إلى عوامل
أولية. ومن ثم يكون كل من ${\OLId{latin-ordinary}{a}}$ و${\OLId{latin-ordinary}{b}}$ غير نسبي، ولدينا
\[
{\OLId{latin-ordinary}{a}}^{\OLId{latin-ordinary}{b}} = \sqrt{{\OLMathNumeral{3}}}^{\log_{\OLMathNumeral{3}} {\OLMathNumeral{4}}} = {\OLMathNumeral{3}}^{{\OLMathNumeral{1}}/{\OLMathNumeral{2}} \cdot \log_{\OLMathNumeral{3}} {\OLMathNumeral{4}}} = ({\OLMathNumeral{3}}^{\log_{\OLMathNumeral{3}}
  {\OLMathNumeral{4}}})^{{\OLMathNumeral{1}}/{\OLMathNumeral{2}}} = {\OLMathNumeral{4}}^{{\OLMathNumeral{1}}/{\OLMathNumeral{2}}}= {\OLMathNumeral{2}},
\]
لأن ${\OLMathNumeral{3}}^{\log_{\OLMathNumeral{3}} {\OLId{latin-ordinary}{x}}} = {\OLId{latin-ordinary}{x}}$.

فالمنطق الحدسي يضبط استدلالًا لا يجيز الخطوة التي قام عليها البرهان الأول.
وإثبات وجود~${\OLId{latin-ordinary}{x}}$ يحقق~$!A({\OLId{latin-ordinary}{x}})$ يقتضي فيه تعيين~${\OLId{latin-ordinary}{x}}$ وإقامة البرهان على
تحققه بــ$!A$، كما صنعنا في البرهان الثاني. وكذلك لا يثبت فيه أن~$!A$ أو
$!B$ تصدق إلا إذا أمكن إثبات واحد منهما.

وفي العبارة الصورية، يحصل المنطق الحدسي بتقييد نظام !!a{derivation}
الكلاسيكي على وجه مخصوص. وهذه، من حيث الرياضيات المحضة، أنظمة استنباط
صورية؛ إلا أن القصد منها تمثيل ضرب من الاستدلال الرياضي. وقد يحمل هذا
الضرب على مذهب فلسفي، كحدسية براور، أو يختار لأنه أصرح وأقرب إلى العيان،
كما في بنائية بيشوب. ويبقى البحث قائمًا في وفاء الرسم الصوري بالدافع غير
الصوري. وعلى اختلاف المذاهب، يصح أن يدرس المنطق الحدسي بوصفه منطقًا
صوريًا قائمًا بنفسه، وفيه من الخصائص ما يجذب كثيرًا من المناطقة الرياضيين.

\end{document}
